\chapter{操作系统接口}
\label{CH:UNIX}

操作系统的职责是在多个程序之间共享一台计算机，并提供比硬件本身所支持的更有用的服务集。
操作系统对底层硬件进行管理和抽象，这样一来，例如，一个字处理器就无需关心具体使用的是哪种类型的磁盘硬件。
操作系统在多个程序之间共享硬件，使得它们可以（或看起来可以）同时运行。
最后，操作系统还为程序之间的交互提供了受控的方式，以便它们可以共享数据或协同工作。

操作系统通过一个接口向用户程序提供服务。
\index{接口设计}
设计一个好的接口被证明是困难的。 一方面，我们希望接口简单且狭窄，因为这使得正确实现它变得更容易。 另一方面，我们又可能倾向于为应用程序提供许多复杂的功能。
解决这种矛盾的关键在于设计出依赖少数几种机制、并通过组合这些机制来提供高度通用性的接口。

本书以一个具体的操作系统作为实例来说明操作系统的概念。 这个操作系统，xv6，提供了由 Ken Thompson 和 Dennis Ritchie 的 Unix 操作系统~\cite{unix} 引入的基本接口，并且模拟了 Unix 的内部设计。 Unix 提供了一个狭窄的接口，其机制能够很好地组合，展现出惊人的通用性。 这个接口非常成功，以至于现代操作系统——BSD、Linux、macOS、Solaris，甚至在较小程度上包括 Microsoft Windows——都提供了类似 Unix 的接口。 理解 xv6 是理解这些系统以及许多其他系统的一个良好起点。

如图 \ref{fig:os} 所示，xv6 采用了传统的
\indextext{内核}
形式，内核是一个特殊的程序，为正在运行的程序提供服务。
每个正在运行的程序，称为
\indextext{进程}，
拥有包含指令、数据和一个栈的内存。 指令实现了程序的运算。 数据是运算所操作的变量。 栈则用于组织程序的函数调用。 一台给定的计算机通常有许多进程，但只有一个内核。

当一个
进程需要调用内核服务时，它会发起一个
\indextext{系统调用}，
这是操作系统接口中的一个调用。
系统调用会进入内核；
内核执行服务并返回。
因此，一个进程会在
\indextext{用户空间}
和
\indextext{内核空间}
中交替执行。

正如后续章节所详述的，内核利用 CPU\footnote{
本文通常使用术语 \indextext{CPU}（中央处理单元的缩写）来指代执行计算的硬件元素。 其他文档（例如 RISC-V 规范）也可能使用处理器（processor）、核心（core）或硬件线程（hart）来代替 CPU。
}
提供的硬件保护机制来确保每个在用户空间执行的进程只能访问其自身的内存。
内核在实现这些保护所需的硬件特权级别下执行；用户程序则在没有这些特权的情况下执行。
当用户程序调用系统调用时，硬件会提升特权级别，并开始执行内核中预先安排好的函数。

\begin{figure}[t]
\center
\includegraphics[scale=0.5]{fig/os.pdf}
\caption{一个内核和两个用户进程。}
\label{fig:os}
\end{figure}

内核提供的系统调用集合就是用户程序所看到的接口。
xv6 内核提供了传统 Unix 内核通常提供的服务和系统调用的一个子集。
图 \ref{fig:api} 列出了 xv6 的所有系统调用。

本章的剩余部分将概述 xv6 的服务——进程、内存、文件描述符、管道和文件系统——并通过代码片段以及讨论 \indextext{shell}（Unix 的命令行用户界面）如何使用它们来说明这些服务。 Shell 对系统调用的使用展示了它们是如何经过精心设计的。

Shell 是一个普通的程序，它从用户那里读取命令并执行它们。
Shell 是一个用户程序而不是内核的一部分，这一事实说明了系统调用接口的强大之处：Shell 并没有什么特别之处。
这也意味着 Shell 很容易被替换；因此，现代 Unix 系统有多种 Shell 可供选择，每种都有其自己的用户界面和脚本功能。
xv6 的 shell 是对 Unix Bourne shell 本质的一个简单实现。

xv6 shell 的实现可以在 \fileref{user/sh.c} 找到。 （蓝色链接是跳转到相关 xv6 源代码的超链接，地址为 \url{https://github.com/mit-pdos/xv6-riscv/}，具体编号指的是 \href{xv6-src-booklet.pdf}{xv6-src-booklet.pdf} 中的页和行号，如同 Lions 对 UNIX 第 6 版的注释~\cite{lions} 一样。一个好的实践是尝试先在您喜欢的开发环境、github 或 PDF 阅读器中阅读源代码，然后再回到本书。在读完本书时，您应该能够在不查阅本书的情况下理解 xv6 的每一行源代码。）

%% 
%% 	Processes and memory
%% 
\section{进程和内存}

一个 xv6 进程包含用户空间内存（指令、数据和栈）以及内核私有的每个进程的状态。
Xv6 
\indextext{分时共享}
进程：它透明地在等待执行的一组进程之间切换可用的 CPU。
当一个进程没有在执行时，xv6 会保存该进程的 CPU 寄存器，并在下次运行该进程时恢复它们。
内核为每个进程关联了一个进程标识符，或称
\indexcode{PID}。

\begin{figure}[t]
\center
\begin{tabular}{ll}
{\bf 系统调用} & {\bf 描述} \\
\midrule
int fork() & 创建一个进程，返回子进程的 PID。 \\
int exit(int status) & 终止当前进程；status 报告给 wait()。不返回。 \\
int wait(int *status) & 等待一个子进程退出；退出状态存放在 *status；返回子进程 PID。 \\
int kill(int pid) & 终止 PID 为 pid 的进程。成功返回 0，错误返回 -1。 \\
int getpid() & 返回当前进程的 PID。 \\
int pause(int n) & 暂停 n 个时钟滴答。 \\
int exec(char *file, char *argv[]) & 加载一个文件并用参数执行它；仅在出错时返回。 \\
char *sbrk(int n) & 将进程的内存增加 n 个零字节。返回新内存的起始地址。 \\
int open(char *file, int flags) & 打开一个文件；flags 指示读/写；返回一个 fd（文件描述符）。 \\
int write(int fd, char *buf, int n) & 从 buf 向文件描述符 fd 写入 n 个字节；返回 n。 \\
int read(int fd, char *buf, int n) & 读取 n 个字节到 buf；返回读取的字节数；文件结束时返回 0。 \\
int close(int fd) & 释放已打开的文件 fd。 \\
int dup(int fd) & 返回一个新的文件描述符，该描述符指向与 fd 相同的文件。\\
int pipe(int p[]) & 创建一个管道，将读文件描述符和写文件描述符放入 p[0] 和 p[1]。 \\
int chdir(char *dir) & 更改当前目录。 \\
int mkdir(char *dir) & 创建一个新目录。 \\
int mknod(char *file, int, int) & 创建一个设备文件。 \\
int fstat(int fd, struct stat *st) & 将打开文件的信息放入 *st。 \\
int link(char *file1, char *file2) & 为文件 file1 创建另一个名字（file2）。 \\
int unlink(char *file) & 移除一个文件。 \\
\end{tabular}
\caption{Xv6 系统调用。除非另有说明，这些调用成功返回 0，出错返回 -1。}
\label{fig:api}
\end{figure}

一个进程可以使用
\indexcode{fork}
系统调用来创建一个新进程。
\lstinline{fork}
为新进程提供调用进程内存的一个精确副本：\lstinline{fork} 将调用进程的指令、数据和栈复制到新进程的内存中。
\lstinline{fork}
在原始进程和新进程中都会返回。
在原始进程中，\lstinline{fork} 返回新进程的 PID。
在新进程中，\lstinline{fork} 返回零。
原始进程和新进程通常被称为
\indextext{父进程}
和
\indextext{子进程}。

例如，考虑下面这个用 C 编程语言~\cite{kernighan} 编写的程序片段：
\begin{lstlisting}[]
int pid = fork();
if(pid > 0){
  printf("parent: child=%d\n", pid);
  pid = wait((int *) 0);
  printf("child %d is done\n", pid);
} else if(pid == 0){
  printf("child: exiting\n");
  exit(0);
} else {
  printf("fork error\n");
}
\end{lstlisting}
\indexcode{exit}
系统调用导致调用进程停止执行并释放资源，如内存和打开的文件。
Exit 接受一个整数状态参数，按约定 0 表示成功，1 表示失败。
\indexcode{wait}
系统调用返回当前进程的一个已退出（或被杀死）的子进程的 PID，并将子进程的退出状态复制到传递给 wait 的地址；如果调用者的子进程都没有退出，
\indexcode{wait}
会等待其中一个退出。
如果调用者没有子进程，\lstinline{wait} 立即返回 -1。
如果父进程不关心子进程的退出状态，它可以向
\lstinline{wait}
传递一个零地址。

在示例中，输出行
\begin{lstlisting}[]
parent: child=1234
child: exiting
\end{lstlisting}
可能以任意顺序出现（甚至交错出现），具体取决于父进程还是子进程首先执行其
\indexcode{printf}
调用。
在子进程退出后，父进程的
\indexcode{wait}
返回，导致父进程打印
\begin{lstlisting}[]
parent: child 1234 is done
\end{lstlisting}
尽管子进程开始时拥有父进程内存的副本，但父进程和子进程使用各自独立的内存和寄存器：在一个进程中改变变量不会影响另一个进程。例如，当
\lstinline{wait}
的返回值被存入父进程的
\lstinline{pid}
变量时，它不会改变子进程中的
\lstinline{pid}
变量。子进程中的
\lstinline{pid}
值仍将为零。

\indexcode{exec}
系统调用
用从文件系统中存储的文件加载而来的新内存映像替换调用进程的内存。
该文件必须具有特定的格式，该格式规定了文件的哪部分包含指令、哪部分是数据、从哪条指令开始执行等。Xv6
使用 ELF 格式，第 \ref{CH:MEM} 章将对此进行更详细的讨论。
通常该文件是编译程序源代码的结果。
当
\indexcode{exec}
成功时，它不会返回到调用程序；
相反，从文件加载的指令会从 ELF 头中声明的入口点开始执行。
\lstinline{exec}
接受两个参数：包含可执行文件的文件名和一个字符串参数数组。
例如：
\begin{lstlisting}[]
char *argv[3];

argv[0] = "echo";
argv[1] = "hello";
argv[2] = 0;
exec("/bin/echo", argv);
printf("exec error\n");
\end{lstlisting}
这个片段用运行中的程序
\lstinline{/bin/echo}
的一个实例来替换调用程序，并使用参数列表
\lstinline{echo}
\lstinline{hello} 运行它。
大多数程序忽略参数数组的第一个元素，按照惯例该元素是程序的名称。

xv6 shell 代表用户使用上述调用来运行程序。Shell 的主要结构很简单；参见
\lstinline{main} 
\lineref{user/sh.c:/main/}。
主循环通过
\indexcode{getcmd}
从用户读取一行输入。
然后它调用 
\lstinline{fork}， 
这会创建 shell 进程的一个副本。父进程调用
\lstinline{wait}，
而子进程则运行命令。例如，如果用户在 shell 中输入了
“\lstinline{echo hello}”，
\lstinline{runcmd}
将以
“\lstinline{echo hello}”
作为参数被调用。
\lstinline{runcmd} 
\lineref{user/sh.c:/runcmd/}
运行实际的命令。对于
“\lstinline{echo hello}”，
它会调用
\lstinline{exec} 
\lineref{user/sh.c:/exec.ecmd/}。
如果
\lstinline{exec}
成功，那么子进程将执行来自
\lstinline{echo}
的指令，而不是来自
\lstinline{runcmd}
的指令。
在某个时刻，
\lstinline{echo}
会调用
\lstinline{exit}，
这将导致父进程在
\lstinline{main}
函数中的
\lstinline{wait}
调用处返回
\lineref{user/sh.c:/main/}。

您可能会好奇为什么
\indexcode{fork}
和
\indexcode{exec}
不合并成一个单独的调用；我们稍后会看到
shell 在其 I/O 重定向的实现中利用了这种分离。
为了避免创建一个重复进程然后立即替换它（使用 \lstinline{exec}）所造成的浪费，
操作系统内核通过使用诸如写时复制（参见第 \ref{CH:PGFAULTS} 节）之类的虚拟内存技术，针对这种用例优化了
\lstinline{fork}
的实现。

Xv6 隐式地分配了大部分用户空间内存：
\indexcode{fork}
分配了子进程复制父进程内存所需的内存，而
\indexcode{exec}
分配了足够容纳可执行文件的内存。
一个在运行时需要更多内存（例如，为了
\indexcode{malloc}）的进程
可以调用
\lstinline{sbrk(n)}
将其数据内存增加
\lstinline{n}
个零字节；
\indexcode{sbrk}
返回新内存的位置。

%% 
%% 	I/O and File descriptors
%% 
\section{I/O 和文件描述符}

一个
\indextext{文件描述符} 
是一个小整数，代表一个由内核管理的、进程可以从中读取或向其中写入的对象。
一个进程可以通过打开一个文件、目录或设备，或者通过创建一个管道，或者通过复制一个现有的描述符来获得一个文件描述符。
为简单起见，我们通常将文件描述符所指的对象称为“文件”；
文件描述符接口抽象了文件、管道和设备之间的差异，使它们都看起来像字节流。
我们将输入和输出称为 \indextext{I/O}。

在内部，xv6 内核将文件描述符用作进程私有表的一个索引，因此每个进程都有一个从零开始的私有文件描述符空间。
按照惯例，进程从文件描述符 0（标准输入）读取，向文件描述符 1（标准输出）写入，并向文件描述符 2（标准错误）写入错误消息。
正如我们将看到的，shell 利用这一约定来实现 I/O 重定向和管道。Shell 确保它总是打开三个文件描述符
\lineref{user/sh.c:/open..console/}，
默认情况下，这些是控制台的文件描述符。

\lstinline{read}
和
\lstinline{write}
系统调用从由文件描述符指定的已打开文件中读取字节和向其中写入字节。
调用
\lstinline{read(fd},
\lstinline{buf},
\lstinline{n)}
从文件描述符
\lstinline{fd}
读取最多
\lstinline{n}
个字节，
将它们复制到
\lstinline{buf}
中，
并返回实际读取的字节数。
每个指向文件（而非管道等）的文件描述符都关联着一个偏移量。
\lstinline{read}
从当前文件偏移量处读取数据，然后将该偏移量增加所读取的字节数：
随后的一次
\lstinline{read}
将返回第一次
\lstinline{read}
返回的字节之后的字节。
当没有更多字节可读时，
\lstinline{read}
返回零以指示文件结束。

调用
\lstinline{write(fd},
\lstinline{buf},
\lstinline{n)}
将
\lstinline{buf}
中的
\lstinline{n}
个字节写入到文件描述符
\lstinline{fd}
所指向的文件，并返回写入的字节数。
只有在发生错误时，写入的字节数才会少于
\lstinline{n}。
与
\lstinline{read}
类似，
\lstinline{write}
在当前文件偏移量处写入数据，然后将该偏移量增加写入的字节数：
每个
\lstinline{write}
都从上一个停止的地方继续写入。

以下程序片段（构成了程序
\lstinline{cat}
的本质）将数据从标准输入复制到标准输出。如果发生错误，它会将一条消息写入标准错误。
\begin{lstlisting}[]
char buf[512];
int n;

for(;;){
  n = read(0, buf, sizeof buf);
  if(n == 0)
    break;
  if(n < 0){
    fprintf(2, "read error\n");
    exit(1);
  }
  if(write(1, buf, n) != n){
    fprintf(2, "write error\n");
    exit(1);
  }
}
\end{lstlisting}
代码片段中需要注意的重要一点是
\lstinline{cat}
不知道它是从文件、控制台还是管道中读取。
同样地，
\lstinline{cat}
也不知道它是向控制台、文件还是其他地方输出。
文件描述符的使用以及文件描述符 0 是输入、文件描述符 1 是输出的约定，使得
\lstinline{cat}
的实现变得简单。

\lstinline{close}
系统调用
会释放一个文件描述符，使其可以被未来的
\lstinline{open}、
\lstinline{pipe}、
或
\lstinline{dup}
系统调用（见下文）重新使用。
新分配的文件描述符
始终是当前进程中编号最小的未使用描述符。

文件描述符和
\indexcode{fork}
的交互使得 I/O 重定向易于实现。
\lstinline{fork}
将父进程的文件描述符表连同其内存一起复制，这样子进程开始时拥有与父进程完全相同的已打开文件。
系统调用
\indexcode{exec}
会替换调用进程的内存，但会保留其文件表。
这种行为允许 shell 通过如下方式实现 \indextext{I/O 重定向}：先 fork，然后在子进程中关闭并重新打开选定的文件描述符，最后调用 \lstinline{exec} 来运行新程序。
以下是一个 shell 为命令
\lstinline{cat}
\lstinline{<}
\lstinline{input.txt} 所运行代码的简化版本：
\begin{lstlisting}[]
char *argv[2];

argv[0] = "cat";
argv[1] = 0;
if(fork() == 0) {
  close(0);
  open("input.txt", O_RDONLY);
  exec("cat", argv);
}
\end{lstlisting}
在子进程关闭文件描述符 0 之后，
\lstinline{open}
保证会使用那个文件描述符
来打开新文件
\lstinline{input.txt}：
0 将是最小可用的文件描述符。
然后，\lstinline{cat}
执行时，它的文件描述符 0（标准输入）指向
\lstinline{input.txt}。
父进程的文件描述符不会受到这个序列的影响，因为它只修改了子进程的描述符。

xv6 shell 中用于 I/O 重定向的代码正是以这种方式工作的
\lineref{user/sh.c:/case.REDIR/}。
回顾一下，在代码的这一点，shell 已经 fork 了子 shell，并且
\lstinline{runcmd}
将调用
\lstinline{exec}
来加载新程序。

\lstinline{open}
的第二个参数由一组标志组成，以位的形式表示，用于控制 \lstinline{open}
的行为。这些可能的值在文件控制（fcntl）头文件中定义
\linerefs{kernel/fcntl.h:/RDONLY/,/TRUNC/}：
\lstinline{O_RDONLY}、
\lstinline{O_WRONLY}、
\lstinline{O_RDWR}、
\lstinline{O_CREATE} 和
\lstinline{O_TRUNC}，
它们指示 \lstinline{open} 分别以只读、只写、读写方式打开文件，如果文件不存在则创建它，以及将文件截断为零长度。

现在应该清楚为什么将
\lstinline{fork}
和 
\lstinline{exec} 
分开调用是有帮助的了：在两者之间，shell 有机会重定向子进程的 I/O，而不会干扰主 shell 自身的 I/O 设置。可以想象一个假设的组合
\lstinline{forkexec} 系统调用，但是使用这样一个调用来执行 I/O 重定向的选项似乎很笨拙。
Shell 可以在调用 \lstinline{forkexec} 之前修改自身的 I/O 设置（然后再撤销这些修改）；或者
\lstinline{forkexec} 可以接受用于指示 I/O 重定向的参数；
或者（最不吸引人的）每个像 \lstinline{cat} 这样的程序都必须被教会如何执行自身的 I/O 重定向。

尽管
\lstinline{fork}
复制了文件描述符表，但每个底层的文件偏移量在父进程和子进程之间是共享的。
考虑这个例子：
\begin{lstlisting}[]
if(fork() == 0) {
  write(1, "hello ", 6);
  exit(0);
} else {
  wait(0);
  write(1, "world\n", 6);
}
\end{lstlisting}
在这个片段的最后，附加到文件描述符 1 的文件将包含数据
\lstinline{hello}
\lstinline{world}。
父进程中的
\lstinline{write}
（由于
\lstinline{wait}
的存在，它会在子进程结束后运行）
会从子进程的
\lstinline{write}
结束的地方继续写入。
这种行为有助于从一系列的 shell 命令（例如
\lstinline{(echo}
\lstinline{hello};
\lstinline{echo}
\lstinline{world)}
\lstinline{>output.txt}）中产生顺序输出。

\lstinline{dup}
系统调用复制一个现有的文件描述符，返回一个指向相同底层 I/O 对象的新文件描述符。
两个文件描述符共享一个偏移量，就像通过
\lstinline{fork}
复制的文件描述符一样。
这是将
\lstinline{hello}
\lstinline{world}
写入文件的另一种方式：
\begin{lstlisting}[]
fd = dup(1);
write(1, "hello ", 6);
write(fd, "world\n", 6);
\end{lstlisting}

如果两个文件描述符是通过一系列
\lstinline{fork}
和
\lstinline{dup}
调用从同一个原始文件描述符派生出来的，那么它们共享一个偏移量。
否则，即使它们是通过对同一个文件的
\lstinline{open}
调用得到的，也不会共享偏移量。
\lstinline{dup} 
允许 shell 实现如下命令：
\lstinline{ls}
\lstinline{existing-file}
\lstinline{non-existing-file}
\lstinline{>}
\lstinline{tmp1}
\lstinline{2>&1}。
其中的
\lstinline{2>&1}
告诉 shell 为命令提供一个文件描述符 2，该描述符是描述符 1 的一个副本。
那么，已存在文件的名称和不存在文件的错误消息都将出现在文件
\lstinline{tmp1}
中。
xv6 shell 不支持对错误文件描述符的 I/O 重定向，但是现在您知道如何实现它了。

文件描述符是一个强大的抽象，
因为它们隐藏了它们所连接的具体对象的细节：
一个向文件描述符 1 写入的进程可能正在写入一个文件、一个像控制台这样的设备，或者一个管道。
%% 
%% 	Pipes
%% 
\section{管道}

一个
\indextext{管道} 
是一个小的内核缓冲区，它以一对文件描述符的形式暴露给进程，一个用于读取，一个用于写入。
向管道的一端写入数据，会使这些数据可以从管道的另一端读取。
管道为进程提供了一种通信方式。

下面的示例代码运行程序
\lstinline{wc}，
并将其标准输入连接到一个管道的读取端。
\begin{lstlisting}[]
int p[2];
char *argv[2];

argv[0] = "wc";
argv[1] = 0;

pipe(p);
if(fork() == 0) {
  close(0);
  dup(p[0]);
  close(p[0]);
  close(p[1]);
  exec("/bin/wc", argv);
} else {
  close(p[0]);
  write(p[1], "hello world\n", 12);
  close(p[1]);
}
\end{lstlisting}
程序调用
\lstinline{pipe}，
它会创建一个新管道并将读和写文件描述符记录在数组
\lstinline{p}
中。
在
\lstinline{fork}
之后，父进程和子进程都有指向该管道的文件描述符。
子进程调用 \lstinline{close} 和 \lstinline{dup} 使文件描述符
零指向管道的读取端，
关闭
\lstinline{p}
中的文件描述符，
然后调用 \lstinline{exec} 来运行
\lstinline{wc}。
当 
\lstinline{wc}
从其标准输入读取时，它将从管道中读取。
父进程关闭管道的读取端，
向管道写入数据，
然后关闭写入端。

如果没有数据可用，对管道的
\lstinline{read}
会等待，直到有数据被写入，或者所有指向管道写入端的文件描述符都被关闭；
在后一种情况下，
\lstinline{read}
将返回 0，就像到达了数据文件的末尾一样。
\lstinline{read}
会阻塞直到不可能有新数据到达，这一事实是上述例子中子进程必须
在运行
\lstinline{wc}
之前关闭管道的写入端的重要原因之一：如果
\lstinline{wc} 的
某个文件描述符指向了管道的写入端，
\lstinline{wc}
将永远不会看到文件结束符。

xv6 shell 实现了像
\lstinline{grep fork sh.c | wc -l}
这样的管道，其方式与上述代码类似
\lineref{user/sh.c:/case.PIPE/}。
子进程创建一个管道来连接管道的左端和右端。然后它调用
\lstinline{fork}
和
\lstinline{runcmd}
来处理管道的左端，
并再次调用
\lstinline{fork}
和
\lstinline{runcmd}
来处理管道的右端，然后等待两者都完成。
管道的右端本身可能是一个包含管道的命令（例如，
\lstinline{a}
\lstinline{|}
\lstinline{b}
\lstinline{|}
\lstinline{c})，
这又会 fork 出两个新的子进程（一个用于
\lstinline{b}
，一个用于
\lstinline{c}）。
因此，shell 可能会创建一棵进程树。这棵树的叶子节点是命令，内部节点则是等待左右子节点完成的进程。

%% 
%% 	File system
%% 
\section{文件系统}

xv6 文件系统提供数据文件，其中包含未解释的字节数组，以及目录，其中包含指向数据文件和其他目录的具名引用。
目录构成一棵树，从一个称为
\indextext{根目录}
的特殊目录开始。
像
\lstinline{/a/b/c}
这样的
\indextext{路径}
指的是根目录
\lstinline{/}
下，名为
\lstinline{a}
的目录内，名为
\lstinline{b}
的目录内，名为
\lstinline{c}
的文件或目录。
不以
\lstinline{/}
开头的路径相对于调用进程的
\indextext{当前目录}
进行求值，当前目录可以通过
\lstinline{chdir}
系统调用进行更改。
下面两个代码片段打开的是同一个文件（假设所有涉及的目录都存在）：
\begin{lstlisting}[]
chdir("/a");
chdir("b");
open("c", O_RDONLY);

open("/a/b/c", O_RDONLY);
\end{lstlisting}
第一个片段将进程的当前目录更改为
\lstinline{/a/b}；
第二个片段既不引用也不更改进程的当前目录。

有一些系统调用可以创建新文件和目录：
\lstinline{mkdir}
创建一个新目录，
使用带
\lstinline{O_CREATE}
标志的 \lstinline{open}
创建一个新的数据文件，
而
\lstinline{mknod}
创建一个新的设备文件。
下面的例子展示了这三种操作：
\begin{lstlisting}[]
mkdir("/dir");
fd = open("/dir/file", O_CREATE|O_WRONLY);
close(fd);
mknod("/console", 1, 1);
\end{lstlisting}
\lstinline{mknod}
创建一个指向设备的特殊文件。
与一个设备文件关联的是主设备号和次设备号
（\lstinline{mknod}
的两个参数），
它们唯一地标识一个内核设备。
当进程随后打开一个设备文件时，内核会将
\lstinline{read}
和
\lstinline{write}
系统调用转向内核设备实现，而不是将它们传递给文件系统。

文件的名字和文件本身是不同的；
同一个底层文件，称为
\indextext{inode}，
可以有多个名字，
称为
\indextext{链接}。
每个链接由目录中的一个条目组成；
该条目包含一个文件名和一个指向 inode 的引用。
一个 inode 保存着关于一个文件的
\indextext{元数据}，
包括其类型（文件、目录或设备）、
其长度、
文件内容在磁盘上的位置，
以及指向该文件的链接数。

\lstinline{fstat}
系统调用
从文件描述符所引用的 inode 中检索信息。
它填充一个
\lstinline{struct}
\lstinline{stat}，
该结构在
\lstinline{stat.h} \fileref{kernel/stat.h} 中定义如下：
\begin{lstlisting}[]
#define T_DIR     1   // 目录
#define T_FILE    2   // 文件
#define T_DEVICE  3   // 设备

struct stat {
  int dev;     // 文件系统的磁盘设备
  uint ino;    // Inode 编号
  short type;  // 文件类型
  short nlink; // 指向文件的链接数
  uint64 size; // 文件大小（字节）
};
\end{lstlisting}

\lstinline{link}
系统调用创建另一个指向与现有文件相同 inode 的文件系统名字。
这个片段创建了一个新文件，它同时拥有名字
\lstinline{a}
和
\lstinline{b}。
\begin{lstlisting}[]
open("a", O_CREATE|O_WRONLY);
link("a", "b");
\end{lstlisting}
读取或写入
\lstinline{a}
与读取或写入
\lstinline{b}
是相同的。
每个 inode 由一个唯一的
\textit{inode}
\textit{编号}标识。
在上面的代码序列之后，可以通过检查
\lstinline{fstat}
的结果来判断
\lstinline{a}
和
\lstinline{b}
是否指向相同的底层内容：
两者将返回相同的 inode 编号 
(\lstinline{ino})，
并且
\lstinline{nlink}
计数将被设置为 2。

\lstinline{unlink}
系统调用从文件系统中移除一个名字。
只有当文件的链接计数为零并且没有文件描述符指向它时，文件的 inode 和存放其内容的磁盘空间才会被释放。
因此，如果在上一个代码序列后加上
\begin{lstlisting}[]
unlink("a");
\end{lstlisting}
inode 和文件内容仍然可以通过
\lstinline{b}
访问到。
此外，
\begin{lstlisting}[]
fd = open("/tmp/xyz", O_CREATE|O_RDWR);
unlink("/tmp/xyz");
\end{lstlisting}
是一种创建没有名字的临时 inode 的惯用方法，当进程关闭
\lstinline{fd}
或退出时，该临时 inode 会被自动清理掉。

Unix 提供了可以从 shell 调用的文件实用程序，作为用户级程序，例如
\lstinline{mkdir}、
\lstinline{ln}、
和
\lstinline{rm}。
这种设计允许任何人通过添加新的用户级程序来扩展命令行界面。事后看来，这个计划似乎是显而易见的，但在 Unix 时代设计的其他系统通常将这些命令内置于 shell 中（并且将 shell 内置于内核中）。

一个例外是
\lstinline{cd}，
它被内置于 shell 中
\lineref{user/sh.c:/if.cmd.0..==..c./}。
\lstinline{cd}
必须更改 shell 自身的当前工作目录。如果
\lstinline{cd}
作为一个普通命令运行，那么 shell 会 fork 一个子进程，子进程会运行
\lstinline{cd}，
而
\lstinline{cd}
会更改
\textit{子进程}的
工作目录。父进程（即
shell）的工作目录不会改变。
%% 
%% 	Real world
%% 
\section{现实世界}

Unix 将“标准”文件描述符、管道以及方便的 shell 语法（用于操作它们）结合在一起，是编写通用、可重用程序方面的一项重大进步。
这个想法激发了一种“软件工具”文化，这是 Unix 强大和流行的主要原因之一，而 shell 是第一个所谓的“脚本语言”。
Unix 系统调用接口至今仍存在于 BSD、Linux 和 macOS 等系统中。

Unix 系统调用接口已通过可移植操作系统接口（POSIX）标准进行了标准化。
Xv6
\textit{不}
符合 POSIX 标准：它缺少许多系统调用（包括像
\lstinline{lseek}
这样的基本调用），
并且它提供的许多系统调用也与标准不同。
我们对 xv6 的主要目标是简单性和清晰性，同时提供一个简单的类 UNIX 系统调用接口。
有些人通过添加更多的系统调用和一个简单的 C 库来扩展 xv6，以便能够运行基本的 Unix 程序。然而，现代内核提供了比 xv6 多得多的系统调用和更多种类的内核服务。例如，它们支持网络、窗口系统、用户级线程、众多设备的驱动程序等等。现代内核不断快速发展，并提供许多超越 POSIX 的特性。

Unix 通过单一的文件名和文件描述符接口集合，统一了对多种类型资源（文件、目录和设备）的访问。
这个想法可以扩展到更多种类的资源；
一个很好的例子是 Plan 9~\cite{Presotto91plan9}，
它将“资源即文件”的概念应用于网络、图形等。
然而，大多数源自 Unix 的操作系统并没有遵循这条道路。

文件系统和文件描述符是强大的抽象。
即便如此，操作系统接口也有其他模型。
Unix 的前身 Multics 以使其看起来像内存的方式抽象了文件存储，从而产生了一种截然不同的接口风格。
Multics 设计的复杂性直接影响了 Unix 的设计者，他们旨在构建更简单的东西。

Xv6 没有提供用户或保护一个用户免受另一个用户侵害的概念；用 Unix 的术语来说，所有 xv6 进程都以 root 身份运行。

本书探讨了 xv6 如何实现其类 Unix 接口，但这些思想和方法不仅适用于 Unix。
任何操作系统都必须将多路复用进程到底层硬件上，将进程彼此隔离，并提供受控的进程间通信机制。
学习完 xv6 之后，您应该能够审视其他更复杂的操作系统，并在这些系统中也看到 xv6 所蕴含的概念。

%% 
\section{练习}
%%

\begin{enumerate}

\item 编写一个程序，使用 Unix 系统调用，通过一对管道（每个方向一个）在两个进程之间“乒乓”传递一个字节。测量程序的性能，以每秒交换次数为单位。

\end{enumerate}
